\section{Scaling Pre-Training of ColBERT Models to Nomic Embed Scale}

In order to establish whether pre-training ColBERT models is beneficial over a simple distillation, we propose to pre-train a ColBERT model on the widely-known datasets used by Nomic Embed~\cite{nussbaum2025nomic}. 
This approach has several advantages: first, it is a robust choice of public data that is easily accessible through the Hugging Face dataset library, which also facilitates implementation. Secondly, a ModernBERT dense model has been trained on this mixture, allowing for easy comparison.

We thus train various models to explore the effect of each training phase: unsupervised contrastive, supervised contrastive, and knowledge distillation. All the different settings are illustrated in Figure~\ref{2602:fig:pipelines}.
The baseline is a simple knowledge distillation step performed on top of the ModernBERT-embed model that has undergone both unsupervised and supervised contrastive phases in the dense setting and thus only a KD phase in the ColBERT setting (setting (a)). It is the most usual ColBERT training setup.
This model is compared to ColBERT-Zero, for which both the unsupervised and supervised steps are also performed in the multi-vector setting, i.e, the whole pre-training is run in the multi-vector setting (setting (c)). This allows for a fair comparison as both models use the same original base model and the exact same data, the only difference being whether the contrastive phases are performed in the dense or late interaction setting.
As the authors also released the dense model obtained after the unsupervised phase only, we study additionally the impact of skipping the unsupervised ColBERT pre-training but still performing the supervised contrastive phase in the ColBERT setting before KD, allowing to study whether the large-scale unsupervised phase is required or simply adding the supervised phase is enough (setting (b)). As the unsupervised phase is, by far, the most expensive step (almost ten times more than supervised + KD, see Table~\ref{2602:tab:hyperparams}), being able to leverage the one already performed by the dense model would be very interesting. 

%An illustration of the different training phases can be found in Figure~\ref{2602:fig:training_phase} and one of the three different training pipeline compared can be found in Figure~\ref{2602:fig:pipelines}.

\begin{figure}
  \centering
  \resizebox{\textwidth}{!}{\usetikzlibrary{arrows.meta,positioning,calc,fit,backgrounds,patterns,decorations.pathreplacing}


% Colors from all three diagrams
\definecolor{querycolor}{HTML}{6B9CD2}
\definecolor{doccolor}{HTML}{7EBF7E}
\definecolor{hardnegcolor}{HTML}{E8A0A0}
\definecolor{poscolor}{HTML}{228B22}
\definecolor{negcolor}{HTML}{8B2222}
\definecolor{arrowgray}{HTML}{666666}
\definecolor{losscolor}{HTML}{E67E22}
\definecolor{highscore}{HTML}{2ECC71}
\definecolor{lowscore}{HTML}{E74C3C}
\definecolor{scorecolor}{HTML}{E8A0A0}
\definecolor{teachercolor}{HTML}{F5A623}
\definecolor{studentcolor}{HTML}{9B59B6}
\definecolor{framegray}{HTML}{444444}
\definecolor{midscore}{HTML}{F1C40F}
\definecolor{klcolor}{HTML}{E91E63}

% ================== ALIGNMENT KNOBS ==================
\def\TopY{1.2cm}          % common top line for the 3 main sections
\def\LossTopY{-11.6cm}    % common top line for the orange loss boxes (lower = longer arrows)
\def\LossBoxH{3.5cm}      % make all loss boxes same height so their bottoms align

\def\XsecA{0cm}
\def\XsecB{12.5cm}
\def\XsecC{24cm}
% =====================================================

\begin{tikzpicture}[
    node distance=0.8cm and 3.5cm,
    box/.style={rectangle,rounded corners=8pt,minimum width=3.5cm,minimum height=1.0cm,align=center,font=\small\ttfamily,draw=#1!60,line width=1pt,fill=#1!25},
    query/.style={box=querycolor},
    doc/.style={box=doccolor},
    posdoc/.style={box=doccolor},
    negdoc/.style={box=hardnegcolor},
    posarr/.style={-{Stealth[length=3mm]}, line width=1.5pt, poscolor},
    negarr/.style={-{Stealth[length=2.5mm]}, line width=0.8pt, negcolor, dashed},
    score/.style={rectangle,rounded corners=3pt,minimum width=0.7cm,minimum height=0.5cm,align=center,font=\tiny\sffamily\bfseries,draw=#1!80,line width=1pt,fill=#1!40},
    pairlabel/.style={font=\small\itshape, text=poscolor},
    neglabel/.style={font=\small\itshape, text=negcolor},
    frame/.style={rectangle,rounded corners=12pt,draw=arrowgray!50,line width=1.5pt,inner sep=12pt},
    modelbox/.style={rectangle,rounded corners=8pt,minimum width=2.8cm,minimum height=3.5cm,align=center,font=\normalsize\sffamily\bfseries,draw=arrowgray!60,line width=1pt,fill=#1!25},
    arr/.style={-{Stealth[length=3mm]}, line width=1pt, arrowgray},
    looparr/.style={-{Stealth[length=3mm]}, line width=1.5pt, losscolor, dashed}
]

% Define common vertical position for all loss boxes
\newcommand{\lossboxtop}{-12cm}

% ============================================================
% SECTION 1: IN-BATCH CONTRASTIVE (Unsupervised)
% ============================================================
\begin{scope}[shift={(\XsecA,0)}]
\node[query, anchor=north] (q1) at (0,\TopY) {Query};
\node[query, below=of q1] (q2) {Query};
\node[query, below=of q2] (q3) {Query};
\node[query, below=of q3] (q4) {Query};
\node[doc, right=of q1] (d1) {Document};
\node[doc, right=of q2] (d2) {Document};
\node[doc, right=of q3] (d3) {Document};
\node[doc, right=of q4] (d4) {Document};
\draw[posarr] (q1.east) -- (d1.west);
\draw[posarr] (q2.east) -- (d2.west);
\draw[posarr] (q3.east) -- (d3.west);
\draw[posarr] (q4.east) -- (d4.west);
\draw[negarr] (q1.east) -- (d2.west);
\draw[negarr] (q1.east) -- (d3.west);
\draw[negarr] (q1.east) -- (d4.west);
\draw[negarr] (q2.east) -- (d1.west);
\draw[negarr] (q2.east) -- (d3.west);
\draw[negarr] (q2.east) -- (d4.west);
\draw[negarr] (q3.east) -- (d1.west);
\draw[negarr] (q3.east) -- (d2.west);
\draw[negarr] (q3.east) -- (d4.west);
\draw[negarr] (q4.east) -- (d1.west);
\draw[negarr] (q4.east) -- (d2.west);
\draw[negarr] (q4.east) -- (d3.west);



\begin{scope}[on background layer]
\node[rectangle, rounded corners=12pt, draw=arrowgray!50, line width=1.5pt,fit=(q1)(q4)(d1)(d4), inner sep=15pt] (pairframe1) {};
\end{scope}
\node[font=\normalsize\sffamily\bfseries, arrowgray] at ($(pairframe1.north)+(0,0.4)$) {Unsupervised Contrastive Pre-Training};

% Loss box at fixed vertical position
\coordinate (lossanchor1) at ($(q1.west)!0.5!(d1.east)$);
\node[rectangle, rounded corners=8pt, draw=losscolor!60, line width=1.5pt,
      fill=losscolor!10, minimum width=11cm, minimum height=\LossBoxH,
      anchor=north] (lossbox1) at (lossanchor1 |- 0,\LossTopY) {};
\node[font=\small\sffamily\bfseries, losscolor!80!black] at ($(lossbox1.north)-(0,0.4cm)$) {InfoNCE loss with In-Batch Contrastive Pairs};
\coordinate (simstart1) at ($(lossbox1.center) + (-3.5cm, 0.1cm)$);
\node[font=\scriptsize\sffamily\bfseries, querycolor!80!black] at ($(simstart1)+(-0.3cm,0.15cm)$) {$q_i$};
\node[font=\tiny\sffamily, anchor=south] at ($(simstart1)+(0.4cm,0.45cm)$) {$d_1$};
\node[font=\tiny\sffamily, anchor=south] at ($(simstart1)+(1.2cm,0.45cm)$) {$d_2$};
\node[font=\tiny\sffamily, anchor=south] at ($(simstart1)+(2.0cm,0.45cm)$) {$d_3$};
\node[font=\tiny\sffamily, anchor=south] at ($(simstart1)+(2.8cm,0.45cm)$) {$d_4$};
\node[score=highscore] at ($(simstart1)+(0.4cm,0.15cm)$) {0.9};
\node[score=lowscore] at ($(simstart1)+(1.2cm,0.15cm)$) {0.2};
\node[score=lowscore] at ($(simstart1)+(2.0cm,0.15cm)$) {0.1};
\node[score=lowscore] at ($(simstart1)+(2.8cm,0.15cm)$) {0.3};
\draw[poscolor, line width=1.5pt] ($(simstart1)+(0.1cm,-0.2cm)$) -- ++(0.6cm,0);
\node[font=\tiny\sffamily, poscolor!80!black] at ($(simstart1)+(0.4cm,-0.45cm)$) {positive};
\draw[-{Stealth[length=2mm]}, line width=1pt, arrowgray] ($(simstart1)+(3.4cm,0.15cm)$) -- ++(0.8cm,0) node[midway, above, font=\tiny\sffamily] {softmax};
\coordinate (softstart1) at ($(simstart1)+(4.5cm,0cm)$);
\node[font=\tiny\sffamily, anchor=south] at ($(softstart1)+(0.3cm,0.7cm)$) {$d_1$};
\node[font=\tiny\sffamily, anchor=south] at ($(softstart1)+(0.9cm,0.7cm)$) {$d_2$};
\node[font=\tiny\sffamily, anchor=south] at ($(softstart1)+(1.5cm,0.7cm)$) {$d_3$};
\node[font=\tiny\sffamily, anchor=south] at ($(softstart1)+(2.1cm,0.7cm)$) {$d_4$};
\fill[highscore!70] ($(softstart1)+(0.05cm,-0.3cm)$) rectangle ++(0.5cm,0.85cm);
\fill[lowscore!50] ($(softstart1)+(0.65cm,-0.3cm)$) rectangle ++(0.5cm,0.18cm);
\fill[lowscore!50] ($(softstart1)+(1.25cm,-0.3cm)$) rectangle ++(0.5cm,0.10cm);
\fill[lowscore!50] ($(softstart1)+(1.85cm,-0.3cm)$) rectangle ++(0.5cm,0.22cm);
\draw[arrowgray!60, line width=0.8pt, rounded corners=3pt] ($(softstart1)+(-0.05cm,-0.35cm)$) rectangle ++(2.5cm,1.05cm);
\node[font=\tiny\sffamily\bfseries, highscore!80!black, align=center] at ($(softstart1)+(1.2cm,-0.7cm)$) {Negatives: random in-batch};
\node[font=\scriptsize\sffamily\bfseries, losscolor!80!black, inner sep=3pt] at ($(lossbox1.south)+(0,0.5cm)$) {$\mathcal{L} = -\frac{1}{N}\sum_{i=1}^{N} \log \frac{\exp(s(q_i, d_i)/\tau)}{\sum_{j=1}^{N} \exp(s(q_i, d_j)/\tau)}$};
\draw[-{Stealth[length=3mm]}, line width=1.5pt, losscolor] (pairframe1.south) -- (lossbox1.north) node[midway, right, font=\scriptsize\sffamily, losscolor!80!black] {};
\end{scope}

% ============================================================
% SECTION 2: HARD NEGATIVE CONTRASTIVE (Supervised)
% ============================================================
\begin{scope}[shift={(\XsecB,0)}]
\node[query, anchor=north] (hq1) at (0,\TopY) {Query};
\node[posdoc, right=of hq1] (hd1) {Positive Document};
\node[negdoc, below=of hd1] (hd1neg1) {Hard Negative Document};
\node[negdoc, below=of hd1neg1] (hd1neg2) {Hard Negative Document};
\node[query, below=4.4cm of hq1] (hq2) {Query};
\node[posdoc, right=of hq2] (hd2) {Positive Document};
\node[negdoc, below=of hd2] (hd2neg1) {Hard Negative Document};
\node[negdoc, below=of hd2neg1] (hd2neg2) {Hard Negative Document};
\draw[posarr] (hq1.east) -- (hd1.west);
\draw[posarr] (hq2.east) -- (hd2.west);
\draw[negarr] (hq1.east) -- (hd1neg1.west);
\draw[negarr] (hq1.east) -- (hd1neg2.west);
\draw[negarr] (hq1.east) -- (hd2.west);
\draw[negarr] (hq1.east) -- (hd2neg1.west);
\draw[negarr] (hq1.east) -- (hd2neg2.west);
\draw[negarr] (hq2.east) -- (hd1.west);
\draw[negarr] (hq2.east) -- (hd1neg1.west);
\draw[negarr] (hq2.east) -- (hd1neg2.west);
\draw[negarr] (hq2.east) -- (hd2neg1.west);
\draw[negarr] (hq2.east) -- (hd2neg2.west);



\begin{scope}[on background layer]
\node[rectangle, rounded corners=12pt, draw=arrowgray!50, line width=1.5pt,fit=(hq1)(hq2)(hd1)(hd2neg2), inner sep=15pt] (pairframe2) {};
\end{scope}
\coordinate (sharedTitleY) at ($(pairframe2.north)+(0,0.4cm)$);
\node[font=\normalsize\sffamily\bfseries, arrowgray] at (sharedTitleY) {Supervised Contrastive Fine-Tuning};

% Loss box at fixed vertical position (same as section 1)
\coordinate (lossanchor2) at ($(hq1.west)!0.5!(hd1.east)$);
\node[rectangle, rounded corners=8pt, draw=losscolor!60, line width=1.5pt,
      fill=losscolor!10, minimum width=11cm, minimum height=\LossBoxH,
      anchor=north] (lossbox2) at (lossanchor2 |- 0,\LossTopY) {};
\node[font=\small\sffamily\bfseries, losscolor!80!black] at ($(lossbox2.north)-(0,0.4cm)$) {InfoNCE loss with Hard Negatives Contrastive Pairs};
\coordinate (simstart2) at ($(lossbox2.center) + (-3.7cm, -0.1cm)$);
\node[font=\scriptsize\sffamily\bfseries, querycolor!80!black] at ($(simstart2)+(-0.3cm,0.15cm)$) {$q_i$};
\node[font=\tiny\sffamily, anchor=south] at ($(simstart2)+(0.4cm,0.45cm)$) {$d_i^+$};
\node[font=\tiny\sffamily, anchor=south, hardnegcolor!80!black] at ($(simstart2)+(1.1cm,0.45cm)$) {$h_1$};
\node[font=\tiny\sffamily, anchor=south, hardnegcolor!80!black] at ($(simstart2)+(1.8cm,0.45cm)$) {$h_2$};
\node[font=\tiny\sffamily, anchor=south] at ($(simstart2)+(2.5cm,0.45cm)$) {$d_j^+$};
\node[font=\tiny\sffamily, anchor=south, hardnegcolor!80!black] at ($(simstart2)+(3.2cm,0.45cm)$) {$h_3$};
\node[font=\tiny\sffamily, anchor=south, hardnegcolor!80!black] at ($(simstart2)+(3.9cm,0.45cm)$) {$h_4$};
\node[score=highscore] at ($(simstart2)+(0.4cm,0.15cm)$) {0.9};
\node[score=lowscore] at ($(simstart2)+(1.1cm,0.15cm)$) {0.7};
\node[score=lowscore] at ($(simstart2)+(1.8cm,0.15cm)$) {0.6};
\node[score=lowscore] at ($(simstart2)+(2.5cm,0.15cm)$) {0.3};
\node[score=lowscore] at ($(simstart2)+(3.2cm,0.15cm)$) {0.5};
\node[score=lowscore] at ($(simstart2)+(3.9cm,0.15cm)$) {0.4};
\draw[poscolor, line width=1.5pt] ($(simstart2)+(0.1cm,-0.2cm)$) -- ++(0.6cm,0);
\node[font=\tiny\sffamily, poscolor!80!black] at ($(simstart2)+(0.4cm,-0.45cm)$) {positive};
\draw[hardnegcolor!80, line width=1pt, decorate, decoration={brace, amplitude=3pt, mirror}] ($(simstart2)+(0.85cm,-0.2cm)$) -- ($(simstart2)+(2.05cm,-0.2cm)$);
\node[font=\tiny\sffamily, hardnegcolor!80!black] at ($(simstart2)+(1.45cm,-0.5cm)$) {hard neg};
\draw[-{Stealth[length=2mm]}, line width=1pt, arrowgray] ($(simstart2)+(4.5cm,0.15cm)$) -- ++(0.8cm,0) node[midway, above, font=\tiny\sffamily] {softmax};
\coordinate (softstart2) at ($(simstart2)+(5.6cm,0.15cm)$);
\node[font=\tiny\sffamily, anchor=south] at ($(softstart2)+(0.25cm,0.7cm)$) {$d_i^+$};
\node[font=\tiny\sffamily, anchor=south] at ($(softstart2)+(0.75cm,0.7cm)$) {$h_1$};
\node[font=\tiny\sffamily, anchor=south] at ($(softstart2)+(1.25cm,0.7cm)$) {$h_2$};
\node[font=\tiny\sffamily, anchor=south] at ($(softstart2)+(1.75cm,0.7cm)$) {...};
\fill[highscore!70] ($(softstart2)+(0.05cm,-0.3cm)$) rectangle ++(0.4cm,0.75cm);
\fill[hardnegcolor!70] ($(softstart2)+(0.55cm,-0.3cm)$) rectangle ++(0.4cm,0.45cm);
\fill[hardnegcolor!70] ($(softstart2)+(1.05cm,-0.3cm)$) rectangle ++(0.4cm,0.35cm);
\fill[lowscore!50] ($(softstart2)+(1.55cm,-0.3cm)$) rectangle ++(0.4cm,0.15cm);
\draw[arrowgray!60, line width=0.8pt, rounded corners=3pt] ($(softstart2)+(-0.05cm,-0.35cm)$) rectangle ++(2.1cm,1.05cm);
\node[font=\tiny\sffamily\bfseries, highscore!80!black, align=center] at ($(softstart2)+(1.0cm,-0.6cm)$) {Negatives: mined hard negatives};
\node[font=\scriptsize\sffamily\bfseries, losscolor!80!black, inner sep=3pt] at ($(lossbox2.south)+(0,0.5cm)$) {$\mathcal{L} = -\frac{1}{N}\sum_{i=1}^{N} \log \frac{\exp(s(q_i, d_i^+)/\tau)}{\exp(s(q_i, d_i^+)/\tau) + \sum_{h \in H_i} \exp(s(q_i, h)/\tau) + \sum_{j \neq i} \exp(s(q_i, d_j)/\tau)}$};
\draw[-{Stealth[length=3mm]}, line width=1.5pt, losscolor] (lossbox2.north |- pairframe2.south) -- (lossbox2.north) node[midway, right, font=\scriptsize\sffamily, losscolor!80!black] {};
\end{scope}

% ============================================================
% SECTION 3: KL DIVERGENCE DISTILLATION
% ============================================================
\begin{scope}[shift={(\XsecC,0)}]

% Redefine styles for this section with smaller sizes
\tikzset{
    smallbox/.style={rectangle,rounded corners=8pt,minimum width=2.4cm,minimum height=0.8cm,align=center,font=\small\sffamily,draw=#1!60,line width=1pt,fill=#1!25},
    smallquery/.style={smallbox=querycolor},
    smalldoc/.style={smallbox=doccolor},
    smallscore/.style={rectangle,rounded corners=4pt,minimum width=0.8cm,minimum height=0.5cm,align=center,font=\scriptsize\sffamily\bfseries,draw=#1!80,line width=1pt,fill=#1!40},
    smallmodelbox/.style={rectangle,rounded corners=8pt,minimum width=2.4cm,minimum height=1.2cm,align=center,font=\normalsize\sffamily\bfseries,draw=#1!60,line width=1.5pt,fill=#1!30}
}

% ==================== TOP: INPUT ====================
\coordinate (kinputcenter) at (2.8cm,0cm);
\node[smallquery, anchor=north] (kq) at (kinputcenter |- 0,\TopY) {Query};
\node[smalldoc, below=0.3cm of kq] (kd1) {Document 1};
\node[smalldoc, below=0.25cm of kd1] (kd2) {Document 2};
\node[smalldoc, below=0.25cm of kd2] (kd3) {Document 3};

\begin{scope}[on background layer]
\node[rectangle, rounded corners=12pt, draw=arrowgray!50, line width=1.5pt, fit=(kq)(kd3), inner sep=15pt] (kinputframe) {};
\end{scope}

% Title centered above input
\node[font=\normalsize\sffamily\bfseries, arrowgray] at (kinputframe.north |- sharedTitleY) {KL Divergence Distillation};

% ==================== MIDDLE: TEACHER & STUDENT ====================
\node[smallmodelbox=teachercolor] (kteacher) at ($(kinputframe.south)+(-2.0cm,-1.8cm)$) {Gemma\\(Teacher)};
\node[smallmodelbox=studentcolor] (kstudent) at ($(kinputframe.south)+(2.0cm,-1.8cm)$) {ColBERT\\(Student)};

% ==================== TEACHER OUTPUTS ====================
\node[smalldoc, minimum width=2cm] (ktd1) at ($(kteacher.south)+(-0.5,-1.5cm)$) {Doc 1};
\node[smalldoc, below=0.2cm of ktd1, minimum width=2cm] (ktd2) {Doc 2};
\node[smalldoc, below=0.2cm of ktd2, minimum width=2cm] (ktd3) {Doc 3};

\node[smallscore=highscore, right=0.2cm of ktd1] (kts1) {0.85};
\node[smallscore=midscore, right=0.2cm of ktd2] (kts2) {0.42};
\node[smallscore=lowscore, right=0.2cm of ktd3] (kts3) {0.12};

\begin{scope}[on background layer]
\node[rectangle, rounded corners=12pt, draw=teachercolor!50, line width=1.5pt, fit=(ktd1)(ktd3)(kts1)(kts3), inner sep=10pt] (ktoutframe) {};
\end{scope}

% ==================== STUDENT OUTPUTS ====================
\node[smalldoc, minimum width=2cm] (ksd1) at ($(kstudent.south)+(-0.5,-1.5cm)$) {Doc 1};
\node[smalldoc, below=0.2cm of ksd1, minimum width=2cm] (ksd2) {Doc 2};
\node[smalldoc, below=0.2cm of ksd2, minimum width=2cm] (ksd3) {Doc 3};

\node[smallscore=highscore, right=0.2cm of ksd1] (kss1) {0.71};
\node[smallscore=midscore, right=0.2cm of ksd2] (kss2) {0.38};
\node[smallscore=lowscore, right=0.2cm of ksd3] (kss3) {0.19};

\begin{scope}[on background layer]
\node[rectangle, rounded corners=12pt, draw=studentcolor!50, line width=1.5pt, fit=(ksd1)(ksd3)(kss1)(kss3), inner sep=10pt] (ksoutframe) {};
\end{scope}

% ==================== BOTTOM: KL LOSS (aligned with other sections) ====================
\node[rectangle, rounded corners=10pt, draw=losscolor!60, line width=1.5pt,
      fill=losscolor!10, minimum width=8cm, minimum height=\LossBoxH,
      anchor=north] (klossbox) at (kinputcenter |- 0,\LossTopY) {};

\node[font=\small\sffamily\bfseries, losscolor!80!black] at ($(klossbox.north)-(0,0.55cm)$) 
     {KL Divergence Loss};

% Distribution bars
\coordinate (ktbar) at ($(klossbox.center) + (-2.2cm, -0.1cm)$);
\node[font=\tiny\sffamily, anchor=east] at ($(ktbar)+(0,0.35cm)$) {D1};
\node[font=\tiny\sffamily, anchor=east] at ($(ktbar)+(0,0cm)$) {D2};
\node[font=\tiny\sffamily, anchor=east] at ($(ktbar)+(0,-0.35cm)$) {D3};

\fill[highscore!70] ($(ktbar)+(0.05cm,0.2cm)$) rectangle ++(1.7cm,0.22cm);
\fill[midscore!70] ($(ktbar)+(0.05cm,-0.12cm)$) rectangle ++(0.84cm,0.22cm);
\fill[lowscore!70] ($(ktbar)+(0.05cm,-0.45cm)$) rectangle ++(0.24cm,0.22cm);
\node[font=\tiny\sffamily\bfseries, teachercolor!80!black] at ($(ktbar)+(0.9cm,0.65cm)$) {Teacher $P$};

\coordinate (ksbar) at ($(klossbox.center) + (0.9cm, -0.1cm)$);
\fill[highscore!70] ($(ksbar)+(0.05cm,0.2cm)$) rectangle ++(1.42cm,0.22cm);
\fill[midscore!70] ($(ksbar)+(0.05cm,-0.12cm)$) rectangle ++(0.76cm,0.22cm);
\fill[lowscore!70] ($(ksbar)+(0.05cm,-0.45cm)$) rectangle ++(0.38cm,0.22cm);

\draw[highscore!80!black, line width=1pt, dashed] ($(ksbar)+(0.05cm,0.2cm)$) rectangle ++(1.7cm,0.22cm);
\draw[midscore!80!black, line width=1pt, dashed] ($(ksbar)+(0.05cm,-0.12cm)$) rectangle ++(0.84cm,0.22cm);
\draw[lowscore!80!black, line width=1pt, dashed] ($(ksbar)+(0.05cm,-0.45cm)$) rectangle ++(0.24cm,0.22cm);
\node[font=\tiny\sffamily\bfseries, studentcolor!80!black] at ($(ksbar)+(0.9cm,0.65cm)$) {Student $Q$};

\draw[-{Stealth[length=2mm]}, line width=1pt, arrowgray] 
    ($(ktbar)+(2cm,0cm)$) -- ($(ksbar)+(-0.1cm,0cm)$);

% Loss formula
\node[font=\scriptsize\sffamily\bfseries, losscolor!80!black] 
    at ($(klossbox.south)+(0,0.5cm)$) {$\mathcal{L}_{KL} = \sum P(d|q) \log \frac{P(d|q)}{Q(d|q)}$};

% ==================== ARROWS ====================
% Input to models
\draw[arr] (kinputframe.south) -- ++(0,-0.4cm) -| (kteacher.north);
\draw[arr] (kinputframe.south) -- ++(0,-0.4cm) -| (kstudent.north);

% Models to outputs with labels
\draw[arr] (kteacher.south) -- (ktoutframe.north)
    node[pos=0.5, left=0.1cm, font=\scriptsize\sffamily\bfseries, teachercolor!80!black] {Teacher}
    node[pos=0.5, right=0.1cm, font=\scriptsize\sffamily\bfseries, teachercolor!80!black] {$P(d|q)$};

\draw[arr] (kstudent.south) -- (ksoutframe.north)
    node[pos=0.5, left=0.1cm, font=\scriptsize\sffamily\bfseries, studentcolor!80!black] {Student}
    node[pos=0.5, right=0.1cm, font=\scriptsize\sffamily\bfseries, studentcolor!80!black] {$Q(d|q)$};

% Outputs to loss
\draw[arr, losscolor] (ktoutframe.south) -- (ktoutframe.south |- klossbox.north);
\draw[arr, losscolor] (ksoutframe.south) -- (ksoutframe.south |- klossbox.north);

\end{scope}

% ============================================================
% SHARED LEGEND (bottom left)
% ============================================================
\node[rectangle, rounded corners=5pt, draw=arrowgray!40, fill=white, inner sep=8pt, anchor=north west] (sharedlegend) at ($(lossbox1.south west)+(0,-0.5)$) {
    \begin{tikzpicture}[scale=0.8]
        \draw[posarr] (0,0) -- (0.8,0);
        \node[font=\small\sffamily, anchor=west, poscolor] at (0.95,0) {Positive Pair};
        \draw[negarr] (4.0,0) -- (4.8,0);
        \node[font=\small\sffamily, anchor=west, negcolor] at (4.95,0) {Negative Pair};
    \end{tikzpicture}
};

\end{tikzpicture}
}
  \caption{An illustration of the three training phases. Unsupervised contrastive pre-training is a very large scale training relying exclusively on large batch sizes to use in-batch negatives. Supervised contrastive fine-tuning is a refinement step that leverages mined hard negative to provide a stronger signal. The knowledge distillation step use a strong teacher to scores various documents and use KL divergence to make the student distribution fits the teachers'. }
  \label{2602:fig:training_phase}
\end{figure}

\begin{figure}
  \centering
  \resizebox{\textwidth}{!}{\definecolor{mlmcolor}{HTML}{4A4A4A}
\definecolor{densecolor}{HTML}{5B9BD5}
\definecolor{colbertcolor}{HTML}{ED7D31}
\definecolor{arrowgray}{HTML}{666666}
\begin{tikzpicture}[
    node distance=1.8cm and 2.2cm,
    box/.style={
        rectangle,
        rounded corners=3pt,
        minimum width=2.4cm,
        minimum height=0.9cm,
        align=center,
        font=\small,
        draw=#1,
        line width=0.8pt,
        fill=#1!12
    },
    dense/.style={box=densecolor},
    colbert/.style={box=colbertcolor},
    mlm/.style={box=mlmcolor},
    arr/.style={-{Stealth[length=2.5mm]}, line width=0.7pt, arrowgray},
    pipelabel/.style={font=\small\bfseries, anchor=west},
    costlabel/.style={font=\small\itshape, anchor=west, text=arrowgray},
]
% === Starting point ===
\node[mlm] (mlm) {MLM\\Pre-training};
% === Pipeline A: KD only (top) ===
\node[dense, right=of mlm, yshift=2.2cm] (a1) {Unsupervised};
\node[dense, right=of a1] (a2) {Supervised};
\node[colbert, right=of a2] (a3) {KD};
% === Pipeline B: Sup + KD (middle) ===
\node[dense, right=of mlm] (b1) {Unsupervised};
\node[colbert, right=of b1] (b2) {Supervised};
\node[colbert, right=of b2] (b3) {KD};
% === Pipeline C: Full ColBERT (bottom) ===
\node[colbert, right=of mlm, yshift=-2.2cm] (c1) {Unsupervised};
\node[colbert, right=of c1] (c2) {Supervised};
\node[colbert, right=of c2] (c3) {KD};
% === Arrows from MLM ===
\draw[arr] (mlm.east) -- ++(0.4,0) |- (a1.west);
\draw[arr] (mlm.east) -- (b1.west);
\draw[arr] (mlm.east) -- ++(0.4,0) |- (c1.west);
% === Pipeline A arrows ===
\draw[arr] (a1) -- (a2);
\draw[arr] (a2) -- (a3);
% === Pipeline B arrows ===
\draw[arr] (b1) -- (b2);
\draw[arr] (b2) -- (b3);
% === Pipeline C arrows ===
\draw[arr] (c1) -- (c2);
\draw[arr] (c2) -- (c3);
% === Pipeline labels + compute costs (right side) ===
\node[pipelabel] at ($(a3.east)+(0.3,0.12)$) {(a) KD only};
\node[costlabel] at ($(a3.east)+(0.3,-0.22)$) {{\raise.17ex\hbox{$\scriptstyle\sim$}}8 GH200-hours};
\node[pipelabel] at ($(b3.east)+(0.3,0.12)$) {(b) Supervised + KD};
\node[costlabel] at ($(b3.east)+(0.3,-0.22)$) {{\raise.17ex\hbox{$\scriptstyle\sim$}}40 GH200-hours};
\node[pipelabel] at ($(c3.east)+(0.3,0.12)$) {(c) ColBERT-Zero};
\node[costlabel] at ($(c3.east)+(0.3,-0.22)$) {{\raise.17ex\hbox{$\scriptstyle\sim$}}408 GH200-hours};
% === Legend ===
\node[dense, minimum width=1.4cm, minimum height=0.55cm] (leg1) at ($(c1.south)!0.5!(c2.south)+(0,-1.0)$) {};
\node[anchor=west, font=\small] at (leg1.east) {~Dense};
\node[colbert, minimum width=1.4cm, minimum height=0.55cm] (leg2) at ($(leg1.east)+(2.4,0)$) {};
\node[anchor=west, font=\small] at (leg2.east) {~ColBERT};
\end{tikzpicture}}
  \caption{The different training pipelines compared in this work. KD only is the most common setup, while ColBERT-Zero is the most expensive. Supervised + KD is an in-between much cheaper than the full pre-training.}
  \label{2602:fig:pipelines}
\end{figure}



We use standard setups for each phase, training for only one epoch each time.
For the unsupervised pre-training, we use the large dataset from Nomic Embed and rely on in-batch negatives. The InfoNCE loss temperature was designated as a learnable parameter with the objective of ascertaining the optimal value for the dataset, and later fixed during training. 
For the supervised phase, we leverage smaller and higher quality datasets with negatives mined by Nomic, while the optimal temperature was determined similarly to the unsupervised phase.
The knowledge distillation setup follows the one of GTE-ModernColBERT, using the MS-MARCO dataset with mined samples scored by \href{https://huggingface.co/BAAI/bge-reranker-v2-gemma}{bge-reranker-v2-gemma}. All the hyperparameters are summarized in Table~\ref{2602:tab:hyperparams}, including the learning rate sweep ranges used for each phase. We use the NanoBEIR evaluator to choose the best values for temperature and learning rate as running the full BEIR~\cite{thakur2021beir} is expensive.
%The query and document lenghts are the same as the supervised phase. Unsupervised pre-training is done on 8 nodes of 4 GH200, while supervised fine-tuning and knowledge distillation use only one node.
\begin{table*}
\centering
\caption{
  Training hyperparameters of the 3 different phases of learning. \\ 
  The learning rates are swept logarithmically in the given ranges, testing 10 values for each phase.
  The query and document lengths include the tokens of the prompts (see Section~\ref{2602:sec:prompts}).
  GPU hours are measured on a cluster of GH200 GPUs, where nodes have 4 GPU each; unsupervised is done on 8 nodes, while supervised and KD on 1 node.
}
\label{2602:tab:hyperparams}
\begin{tabular}{l c c c c c c}
\toprule
\textbf{Phase} & \textbf{LR sweep range} & \textbf{Batch size} & \textbf{Temp.} & \textbf{Query length} & \textbf{Doc. length} & \textbf{GPU hours}\\
\midrule
% Unsupervised            & \(\num{3e-3}, \num{1e-3}, \num{8e-4}, \num{6e-4}, \num{4e-4}, \num{3e-4}, \num{2e-4}, \num{1e-4}, \num{3e-5}, \num{1e-5}\) & 16384 & 0.2\\
% Supervised              & \(\num{2e-5}, \num{1e-5}, \num{2e-6}, \num{1e-6}, \num{8e-7}, \num{6e-7}, \num{4e-7}, \num{2e-7}, \num{1e-7}, \num{8e-8}\) & 64 & 0.2 \\
% KD  & \(\num{1e-3}, \num{8e-4}, \num{3e-4}, \num{1e-4}, \num{6e-5}, \num{4e-5}, \num{2e-5}, \num{1e-5}, \num{3e-7}, \num{1e-7}\) & \(128^\star\) & n.a.\\
Unsupervised  & \([\num{3e-3}, \num{1e-5}]\) & 16384 & 0.2 & 39 & 187 & 368\\
Supervised    & \([\num{2e-5}, \num{8e-8}]\) & 64 & 0.2 & 39 & 519 & 32\\
KD            & \([\num{1e-3},  \num{1e-7}]\) & \(128^\star\) & n.a. & 39 & 519 & 8\\
\bottomrule
\end{tabular}


{\small \(\star\) Effective batch size achieved with gradient accumulation of 2.}
\end{table*}



We use the PyLate~\cite{chaffin2025pylate} library and leverage the implementation of GradCache~\cite{gao2021scaling} to scale the per-GPU batch size arbitrarily without VRAM constraints (as standard gradient accumulation is not applicable for contrastive learning). We also gather the samples from the other GPUs in a multi-GPU setting to leverage all computed representations and scale batch size without additional memory requirement. These two features combined allows to easily reach batch sizes of 16k, which is required for unsupervised training to get plausible in-batch hard negatives~\cite{chen2020simple}. Finally, each batch is composed of only one data source at a time, to prevent shortcut learning~\cite{nussbaum2025nomic}. We release the checkpoints obtained after each phase (unsupervised, supervised, distilled) for all pipelines, enabling the community to use them as starting points for their own experiments.

% Unsupervised pre-training is done on 8 nodes of 4 GH200, while supervised fine-tuning and knowledge distillation use only one node. The total GH200 hours of each phase are respectively, 368, 32 and 8, highlighting that the unsupervised pre-training is --by far-- the most expensive phase. 


% It is worth noting that, besides the limited efforts put into training of ColBERT models (w.r.t dense), one of the reason of the lack of pre-trained ColBERT was the lack of available resources to easily explore these setups.
% Thankfully, PyLate~\cite{chaffin2025pylate} offers different features that make it suitable for running such large-scale pre-training.
% First, the implementation of GradCache~\cite{gao2021scaling} allows to scale the per-GPU batch size arbitrarily without VRAM constraints (as standard gradient accumulation is not applicable for contrastive learning). Then, the option of gathering the samples from the other GPUs in a multi-GPUs settings allows to leverage all the computed representations to scale batch size without additional memory requirement. These two features combined allows to easily reach batch sizes of 16k, which is required for unsupervised training to get plausible in-batch hard negatives.
% Finally, PyLate allows to easily track in-training evaluations such as NanoBEIR on Weights\&Biases, which makes sweeping and monitoring training easier.
